\documentclass[a5paper,10pt]{article}

% https://github.com/InDevRus/filippov-solutions

% Нумерация страниц
\pagenumbering{gobble}

% Настройка полей
\usepackage{geometry}
\geometry{left=1cm}
\geometry{right=1cm}
\geometry{top=1cm}
\geometry{bottom=1.5cm}

% Вставка таблицы
\usepackage{array}
\usepackage{tabularx}

% Инструменты выравнивания формул
\usepackage{amsmath}
\usepackage{mathtools}

% Диакритика в математике
\usepackage{amsfonts}

% Вычисления размеров на ходу
\usepackage{calc}

% Удобные записи производных
\usepackage[italic=true]{derivative}

% Настройки сносок
\usepackage{enotez}

% Длинные знаки векторов
\usepackage{anyfontsize}
\usepackage[b]{esvect}

% Вставка картинок
\usepackage{graphicx}

% Вставка красной строки
\usepackage{indentfirst}
\setlength{\parindent}{1em}

% Условия в командах
\usepackage{xifthen}

% Луа-скрипты
\usepackage{luacode}

% Настройка команд отдельно в математическом и текстовом режимах
\usepackage{mathcommand}

% Для повторения LaTeX кода
\usepackage{multido}

% Конфигурация отступов формул
\delimitershortfall-1sp
\usepackage{mleftright}
\mleftright

% Растягивание объектов
\usepackage{relsize}

% Обтекание текстом
\usepackage{wrapfig}
\usepackage{subcaption}
\usepackage{floatflt}

% Поддержка ввода кириллицы
\usepackage[english, russian]{babel}

% Настройка корневой позиции для компилятора
\usepackage{import}
\providecommand*{\ProjectRootPath}{..}

\import{\ProjectRootPath/preambles/macros/}{theorems}

\import{\ProjectRootPath/preambles/fonts}{font_setup}

% Знак градуса
\usepackage{siunitx}

\import{\ProjectRootPath/preambles/macros/}{operators}
\import{\ProjectRootPath/preambles/macros/}{redefinitions}

\import{\ProjectRootPath/preambles/fonts}{letters_setup}
\import{\ProjectRootPath/preambles/fonts/adjustments}{custom_superscripts}

\import{\ProjectRootPath/preambles/custom}{formatting}
\import{\ProjectRootPath/preambles/custom}{frames}
\import{\ProjectRootPath/preambles/custom}{list_setup}

% TODO: перевести команды на современный стиль объявления
\input{../preambles/plotting}

\begin{document}

$\ddot{x} = -2\dot{x} - x - \dot{x}^2$.
$\tsys{& \dot{x}\br{t} = y \\& \dot{y}\br{t} = -x - 2y - y^2}$.
$\dfrac {dy} {dx} = \dfrac {-x - 2y - y^2} {y} = \dfrac {-x - 2y + O\br{y^2}} {y}$.

Особой будет лишь точка $M\br{0; 0}$.
Линеаризованное уравнение $\dfrac {dy} {dx} = \dfrac {-x - 2y} {y}$
имеет матрицу
$A = \begin{pmatrix} 0 & 1 \\ -1 & -2 \end{pmatrix}$
и характеристическое уравнение $\lambda^2 + 2\lambda + 1 = 0$, $\lambda_{1} = \lambda_{2} = -1$.
Таким образом, $M\br{0; 0}$ -- вырожденный узел, устойчивый для системы.

Асимптотическую устойчивость также можно показать, используя функцию Ляпунова
$v\br{x; y} = 2x^2 + y^2 + xy$,
$\dfrac {dv} {dt} \bigg| = -x^2 - 3y^2 - xy^2 - 2y^3 = -x^2 - y^2\br{x + 1} - 2y^2\br{y + 1} < 0$ при $0 < \sqrt{x^2 + y^2} < \epsilon_{0} = 1$.

$A - \lambda E = \begin{pmatrix} 1 & 1 \\ -1 & -1 \end{pmatrix} \sim \begin{pmatrix} 1 & 1 \end{pmatrix}$, $v_{1} = \begin{pmatrix} 1 \\ -1 \end{pmatrix}$ и касательной будет прямая $y = -x$.

Траектории построим при помощи изоклин $x = -y^2 - \br{k + 2} y$, представленных на левом графике. График справа содержит решения уравнения без малых возмущений $y' = \dfrac {-x - 2y} {y}$ -- кривые $y + x = C \exp\br{- \dfrac {x} {y + x}}$. Из графика видно, что в окрестности особой точки такие кривые пренебрежимо мало отличаются от настоящих решений.

\begin{figure}[!h]
    \begin{minipage}{.5\textwidth}
        \centering
            \begin{tikzpicture}
            \pgfplotsset{
                Graph/.style = {
                    samples = 200,
                    restrict y to domain = -5 : 5,
                },
                DirectionGraph/.style = {
                    samples = 40,
                    quiver = {scale arrows = 0.05},
                    -stealth,
                },
                DirectionGraph1/.style = {
                    DirectionGraph,
                    quiver = {
                        u = {dot_x(x, y) / sqrt(dot_x(x, y) ^ 2 + dot_y(x, y) ^ 2)},
                        v = {dot_y(x, y) / sqrt(dot_x(x, y) ^ 2 + dot_y(x, y) ^ 2)}
                    }
                },
            }
            \begin{axis}[
                trig format plots=rad,
                width = 7.5cm,
                height = 10cm,
                ylabel=$\dot{x}$,
                ymin = -1.2,
                ymax = 1.2,
                xmin = -1.2,
                xmax = 1.2,
                xtick = {-10, ..., 10},
                ytick = {-10, ..., 10},
                minor tick num = 3,
                declare function = {
                    xi(\k, \x) = -x^2 - (2 + \k) * x;
                    dot_x(\x, \y) = y;
                    dot_y(\x, \y) = -x - 2 * y - y^2;
                }]

                \foreach \C in {-6, -5.5, ..., 3, -4.25, -3.75, -3.25, -2.75, -0.75, -0.25, 0.25, 0.75, 1.25} {
                    \addplot[Graph, cyan, dotted, domain = -1.5 : 1.5]
                    ({xi(\C, x)}, {x});
                    \addplot[DirectionGraph1, samples = 15, domain = -1.2 : -0.15]
                    ({xi(\C, x)}, {x});
                    \addplot[DirectionGraph1, samples = 15, domain = 0.15 : 1.2]
                    ({xi(\C, x)}, {x});
    				\addplot[DirectionGraph1, samples = 8, domain = -1.1 : -0.15]
    				({x}, {0});
    				\addplot[DirectionGraph1, samples = 8, domain = 0.15 : 1.1]
    				({x}, {0});
                }
            \end{axis}
        \end{tikzpicture}
    \end{minipage}
    \begin{minipage}{.5\textwidth}
    	\centering
    	\begin{tikzpicture}
    		\pgfplotsset{
    			Graph/.style = {
    				samples = 200,
    				restrict y to domain = -5 : 5,
    			},
    			DirectionGraph/.style = {
    				opacity = 0.5,
    				samples = 40,
    				quiver = {scale arrows = 0.05},
    				-stealth,
    			},
    			DirectionGraph1/.style = {
    				DirectionGraph,
    				quiver = {
    					u = {dot_x(x, y) / sqrt(dot_x(x, y) ^ 2 + dot_y(x, y) ^ 2)},
    					v = {dot_y(x, y) / sqrt(dot_x(x, y) ^ 2 + dot_y(x, y) ^ 2)}
    				}
    			},
    		}
    		\begin{axis}[
    			trig format plots=rad,
    			width = 7.5cm,
    			height = 10cm,
    			ylabel=$\dot{x}$,
    			ymin = -1.2,
    			ymax = 1.2,
    			xmin = -1.2,
    			xmax = 1.2,
    			xtick = {-10, ..., 10},
    			ytick = {-10, ..., 10},
    			minor tick num = 3,
    			declare function = {
    				xi(\k, \x) = -x^2 - (2 + \k) * x;
    				dot_x(\x, \y) = y;
    				dot_y(\x, \y) = -x - 2 * y - y^2;
    				f(\C, \x) = x * (\C - ln(abs(x)));
    				g(\C, \x) = x * (1 - \C + ln(abs(x)));
    			}]
    			
    			\foreach \C in {-6, -5.5, ..., 3, -4.25, -3.75, -3.25, -2.75, -0.75, -0.25, 0.25, 0.75, 1.25} {
    				\addplot[DirectionGraph1, samples = 15, domain = -1.2 : -0.15]
    				({xi(\C, x)}, {x});
    				\addplot[DirectionGraph1, samples = 15, domain = 0.15 : 1.2]
    				({xi(\C, x)}, {x});
    				\addplot[DirectionGraph1, gray, samples = 8, domain = -1.1 : -0.15]
    				{0};
    				\addplot[DirectionGraph1, gray, samples = 8, domain = 0.15 : 1.1]
    				{0};
    			}
    			
    			\addplot[Graph, samples = 10, dashed, domain = -1.5 : 1.5]
    			{-x};
    			    			
    			\foreach \C in {-2.5, -1.4, -0.8, -0.4, 0, 0.5, 0.65, 0.85, 1, 1.1, 2} {
    				\addplot[Graph, domain = -10 : 0]
    				({f(\C, \x)}, {g(\C, \x)});
    				\addplot[Graph, domain = 0 : 10]
    				({f(\C, \x)}, {g(\C, \x)});
    			}
    			\addplot[Graph, domain = -10 : 0]
    			({f(0.25, \x)}, {g(0.25, \x)});
    			\addplot[Graph, domain = 0 : 1.17]
    			({f(0.25, \x)}, {g(0.25, \x)});
    		\end{axis}
    	\end{tikzpicture}
    \end{minipage}
\end{figure}
Вывод: решения $x\br{t}$ при $t \to +\infty$ стремятся к нулю.

\end{document}
